\documentclass[a4paper,12pt]{article}

\usepackage[T2A]{fontenc}
\usepackage[utf8]{inputenc}
\usepackage[russian]{babel}
\usepackage{amsmath,amssymb,amsthm,mathtools}
\usepackage{helvet}
\renewcommand{\familydefault}{\sfdefault}
\usepackage{icomma}
\usepackage{geometry}
\geometry{left=20mm,right=20mm,top=20mm,bottom=22mm}
\usepackage{microtype}
\usepackage{tikz}
\usetikzlibrary{calc,arrows.meta,positioning,shapes.geometric}
\usepackage{pgfplots}
\pgfplotsset{compat=1.18}
\usepackage{tabularx,booktabs,longtable}
\usepackage{enumitem}
\usepackage{hyperref}
\usepackage{parskip}
\usepackage{fancyhdr}
\usepackage{setspace}
\usepackage{xcolor}

\definecolor{accent}{HTML}{008080}
\definecolor{darkaccent}{HTML}{004D4D}
\definecolor{textgray}{HTML}{333333}
\definecolor{softgray}{HTML}{F2F5F5}

\hypersetup{colorlinks=true,linkcolor=darkaccent,urlcolor=darkaccent}
\color{textgray}
\onehalfspacing
\setlength{\parindent}{0pt}
\setlist[itemize]{leftmargin=1.4em,itemsep=0.25em,topsep=0.3em}
\setlist[enumerate]{leftmargin=1.6em,itemsep=0.45em,topsep=0.4em}
\renewcommand{\arraystretch}{1.25}

\pagestyle{fancy}
\fancyhf{}
\fancyhead[L]{\small\textcolor{darkaccent}{Физика: ускорение и график скорости}}
\fancyhead[R]{\small\textcolor{textgray}{Володя Хачатурян}}
\fancyfoot[C]{\small\thepage}
\setlength{\headheight}{25pt}
\renewcommand{\headrulewidth}{0.4pt}
\renewcommand{\headrule}{\hbox to\headwidth{\color{accent}\leaders\hrule height \headrulewidth\hfill}}

\newcommand{\sectionline}{%
  \par\vspace{-0.5em}%
  \noindent\textcolor{accent}{\rule{\linewidth}{0.7pt}}%
  \par\vspace{0.5em}%
}
\newcommand{\key}[1]{\textcolor{darkaccent}{\textbf{#1}}}

\tikzset{
  flowstart/.style={
    ellipse,
    draw=darkaccent,
    line width=0.8pt,
    align=center,
    minimum width=31mm,
    minimum height=8mm,
    inner sep=3pt,
    font=\small
  },
  flowproc/.style={
    rectangle,
    rounded corners=3pt,
    draw=accent,
    line width=0.8pt,
    align=center,
    text width=70mm,
    minimum height=9mm,
    inner sep=4pt,
    font=\small
  },
  flowcond/.style={
    diamond,
    aspect=2.2,
    draw=accent,
    line width=0.8pt,
    align=center,
    text width=34mm,
    inner sep=1.5pt,
    font=\small
  },
  flowarrow/.style={
    -{Stealth[length=3mm,width=2mm]},
    line width=0.9pt,
    draw=darkaccent
  },
  flowlabel/.style={
    font=\small,
    fill=white,
    inner sep=1.5pt,
    text=textgray
  }
}

\begin{document}

% ============================================================
% ТИТУЛЬНЫЙ ЛИСТ
% ============================================================

\begin{titlepage}
\thispagestyle{empty}
\centering
\vspace*{18mm}

{\Large\textcolor{darkaccent}{\textbf{Чек-лист}}\par}
\vspace{7mm}

{\huge\bfseries Ускорение и график скорости\par}
\vspace{4mm}

{\Large Равноускоренное прямолинейное движение\par}

\vspace{18mm}

{\large Физика \textbullet{} подготовка к ОГЭ\par}

\vspace{10mm}

{\large Володя Хачатурян\par}

\vspace{5mm}

{\large \today\par}

\vfill

{\normalsize Лёвин Артём Александрович\par}
{\normalsize эксклюзивно для Володи Хачатуряна\par}

\vspace{18mm}

\begin{tikzpicture}[remember picture,overlay]
  \draw[accent,line width=1.2pt]
  ($(current page.south west)+(18mm,15mm)$)
  .. controls
  ($(current page.south west)+(65mm,27mm)$)
  and
  ($(current page.south east)+(-65mm,7mm)$)
  ..
  ($(current page.south east)+(-18mm,17mm)$);
\end{tikzpicture}

\end{titlepage}

% ============================================================
% ОСНОВНОЙ ЧЕК-ЛИСТ
% ============================================================

\section*{1. Что нужно знать после занятия}
\sectionline

\begin{itemize}
  \item Закон движения задаёт зависимость координаты от времени:
  \[
    x=x(t).
  \]

  \item Для равномерного движения вдоль оси \(Ox\):
  \[
    x=x_0+v_xt.
  \]

  \item Точка встречи двух тел находится из условия
  \[
    x_1(t)=x_2(t).
  \]
  Сначала определяется время встречи, затем оно подставляется в любой закон движения и находится координата встречи.

  \item Ускорение показывает, как быстро изменяется скорость:
  \[
    a_x=
    \frac{v_{x\,\text{кон}}-v_{x\,\text{нач}}}{\Delta t}
    =
    \frac{\Delta v_x}{\Delta t}.
  \]

  \item Единица ускорения в СИ:
  \[
    \text{м}/\text{с}^2.
  \]

  \item При постоянном ускорении скорость меняется по закону
  \[
    v_x=v_{0x}+a_xt.
  \]

  \item График \(v_x(t)\) при постоянном ускорении является прямой. Его пересечение с вертикальной осью даёт \(v_{0x}\), а угловой коэффициент равен \(a_x\).
\end{itemize}

% ============================================================

\section*{2. Повторение: закон движения и точка встречи}
\sectionline

Для двух тел:
\[
  x_1=x_{01}+v_{1x}t,
  \qquad
  x_2=x_{02}+v_{2x}t.
\]

В момент встречи координаты тел одинаковы:
\[
  x_1=x_2.
\]

\key{Пример.}

Пусть
\[
  x_1=24-3t,
  \qquad
  x_2=-4+4t.
\]

Приравниваем координаты:
\[
  24-3t=-4+4t.
\]

Отсюда
\[
  28=7t,
  \qquad
  t=4\,\text{с}.
\]

Теперь находим координату встречи:
\[
  x=24-3\cdot4=12\,\text{м}.
\]

Проверка по второму закону:
\[
  x=-4+4\cdot4=12\,\text{м}.
\]

\key{Критически важно:}
если в задаче требуется \emph{точка встречи}, конечным ответом является координата \(x\). Время встречи служит промежуточной величиной.

% ============================================================

\section*{3. Ускорение: смысл, знак, размерность}
\sectionline

Ускорение характеризует изменение скорости за единицу времени:
\[
  a_x=\frac{\Delta v_x}{\Delta t}.
\]

Если скорость изменилась от
\[
  10\,\text{м}/\text{с}
\]
до
\[
  20\,\text{м}/\text{с}
\]
за
\[
  5\,\text{с},
\]
то
\[
  a_x=
  \frac{20-10}{5}
  =
  2\,\text{м}/\text{с}^2.
\]

Это означает, что каждую секунду проекция скорости увеличивается на
\[
  2\,\text{м}/\text{с}.
\]

По секундам:
\[
  10\to12\to14\to16\to18\to20.
\]

\subsection*{Как читать знак ускорения}

Сам по себе знак \(a_x\) показывает направление изменения \(v_x\).
Для вывода о разгоне или торможении нужно учитывать одновременно знаки \(v_x\) и \(a_x\).

\begin{center}
\begin{tabularx}{\linewidth}{@{}c c X@{}}
\toprule
\(v_x\) & \(a_x\) & Физический смысл \\
\midrule
\(+\) & \(+\) &
тело движется вдоль \(+Ox\) и разгоняется \\
\(+\) & \(-\) &
тело движется вдоль \(+Ox\) и замедляется \\
\(-\) & \(-\) &
тело движется вдоль \(-Ox\) и разгоняется \\
\(-\) & \(+\) &
тело движется вдоль \(-Ox\) и замедляется \\
\bottomrule
\end{tabularx}
\end{center}

\key{Главная идея:}
при разгоне увеличивается \(|v_x|\), при торможении \(|v_x|\) уменьшается.

% ============================================================

\section*{4. Рабочая формула и четыре типа задач}
\sectionline

Основная формула равноускоренного движения для скорости:
\[
  \boxed{v_x=v_{0x}+a_xt}.
\]

По смыслу:

\[
\text{конечная скорость}
=
\text{начальная скорость}
+
\text{изменение скорости}.
\]

Причём
\[
  \Delta v_x=a_xt.
\]

\subsection*{Тип 1. Найти ускорение}

\[
  a_x=\frac{v_x-v_{0x}}{t}.
\]

Например:
\[
  v_{0x}=10\,\text{м}/\text{с},
  \qquad
  v_x=20\,\text{м}/\text{с},
  \qquad
  t=5\,\text{с}.
\]

Тогда
\[
  a_x=
  \frac{20-10}{5}
  =
  2\,\text{м}/\text{с}^2.
\]

\subsection*{Тип 2. Найти конечную скорость}

Пусть
\[
  v_{0x}=10\,\text{м}/\text{с},
  \qquad
  a_x=3\,\text{м}/\text{с}^2,
  \qquad
  t=10\,\text{с}.
\]

Подставляем:
\[
  v_x=10+3\cdot10.
\]

Получаем:
\[
  v_x=40\,\text{м}/\text{с}.
\]

\subsection*{Тип 3. Найти начальную скорость}

Пусть
\[
  v_x=100\,\text{м}/\text{с},
  \qquad
  a_x=5\,\text{м}/\text{с}^2,
  \qquad
  t=20\,\text{с}.
\]

Записываем:
\[
  100=v_{0x}+5\cdot20.
\]

Следовательно,
\[
  100=v_{0x}+100,
\]
\[
  v_{0x}=0\,\text{м}/\text{с}.
\]

\subsection*{Тип 4. Найти время}

Пусть
\[
  v_x=20\,\text{м}/\text{с},
  \qquad
  v_{0x}=2\,\text{м}/\text{с},
  \qquad
  a_x=5\,\text{м}/\text{с}^2.
\]

Получаем:
\[
  20=2+5t,
\]
\[
  18=5t,
\]
\[
  t=\frac{18}{5}=3{,}6\,\text{с}.
\]

% ============================================================

\section*{5. Алгебраический минимум для физики}
\sectionline

После правильной записи физической формулы задача часто превращается в обычное алгебраическое уравнение.

Держите под рукой два базовых шаблона.

\begin{center}
\begin{tabularx}{0.9\linewidth}{@{}X X@{}}
\toprule
Если \(\dfrac{A}{B}=C\) & Если \(A-B=C\) \\
\midrule
\(A=BC\) & \(A=B+C\) \\
\(B=\dfrac{A}{C}\) & \(B=A-C\) \\
\bottomrule
\end{tabularx}
\end{center}

Также полезно помнить:
\[
  A+B=C
  \quad\Rightarrow\quad
  A=C-B,
  \qquad
  B=C-A.
\]

\key{Рабочий принцип:}
сначала подставьте известные числа в физическую формулу, затем решайте полученное уравнение обычными алгебраическими действиями.

% ============================================================
% БЛОК-СХЕМА
% ============================================================

\clearpage
\thispagestyle{empty}

\begin{center}
{\Large\textcolor{darkaccent}{\textbf{Алгоритм: задача на ускорение и скорость}}}
\end{center}

\vspace{4mm}

\begin{center}
\begin{tikzpicture}[node distance=5mm]

  \node[flowstart] (start) {Начало};

  \node[flowproc,below=of start] (data)
  {Выпишите данные с единицами СИ и определите неизвестную величину};

  \node[flowproc,below=of data] (formula)
  {Запишите рабочую формулу\\
  \(v_x=v_{0x}+a_xt\)
  или
  \(a_x=\dfrac{v_x-v_{0x}}{t}\)};

  \node[flowproc,below=of formula] (subst)
  {Подставьте известные числа, неизвестную оставьте буквой};

  \node[flowcond,below=7mm of subst] (simple)
  {Получилось простое уравнение?};

  \node[flowproc,below=7mm of simple] (solve)
  {Решите уравнение: перенесите слагаемые, затем выполните умножение или деление};

  \node[flowcond,below=7mm of solve] (check)
  {Единицы и знак ответа физически разумны?};

  \node[flowproc,below=7mm of check] (answer)
  {Запишите ответ с единицами измерения};

  \node[flowstart,below=of answer] (finish)
  {Конец};

  \node[
    flowproc,
    right=14mm of simple,
    text width=42mm
  ] (algebra)
  {Сначала упростите числовую часть и вспомните базовый алгебраический шаблон};

  \node[
    flowproc,
    right=14mm of check,
    text width=42mm
  ] (repair)
  {Исправьте знак, вычисление или размерность и проверьте ещё раз};

  \draw[flowarrow] (start) -- (data);
  \draw[flowarrow] (data) -- (formula);
  \draw[flowarrow] (formula) -- (subst);
  \draw[flowarrow] (subst) -- (simple);

  \draw[flowarrow]
  (simple) --
  node[flowlabel,left]{да}
  (solve);

  \draw[flowarrow]
  (simple.east) --
  node[flowlabel,above]{нет}
  (algebra.west);

  \draw[flowarrow]
  (algebra.south) |- (solve.east);

  \draw[flowarrow]
  (solve) -- (check);

  \draw[flowarrow]
  (check) --
  node[flowlabel,left]{да}
  (answer);

  \draw[flowarrow]
  (check.east) --
  node[flowlabel,above]{нет}
  (repair.west);

  \draw[flowarrow]
  (repair.south) |- (answer.east);

  \draw[flowarrow]
  (answer) -- (finish);

\end{tikzpicture}
\end{center}

\clearpage

% ============================================================
% ГРАФИКИ
% ============================================================

\section*{6. График зависимости скорости от времени}
\sectionline

Если
\[
  v_x=v_{0x}+a_xt,
\]
то график \(v_x(t)\) является прямой.

\key{По графику можно определить:}

\begin{itemize}
  \item начальную скорость \(v_{0x}\) --- значение \(v_x\) при \(t=0\);

  \item ускорение:
  \[
    a_x=\frac{\Delta v_x}{\Delta t};
  \]

  \item момент остановки --- момент, при котором
  \[
    v_x=0;
  \]

  \item направление движения по знаку скорости:
  \[
    v_x>0
    \Rightarrow
    \text{движение вдоль }+Ox,
  \]
  \[
    v_x<0
    \Rightarrow
    \text{движение вдоль }-Ox.
  \]
\end{itemize}

\subsection*{Пример 1. Формула \(v_x=1+2t\)}

Из формулы сразу получаем:
\[
  v_{0x}=1\,\text{м}/\text{с},
  \qquad
  a_x=2\,\text{м}/\text{с}^2.
\]

Для построения прямой достаточно двух точек:
\[
  t=0
  \Rightarrow
  v_x=1,
\]
\[
  t=1
  \Rightarrow
  v_x=3.
\]

\begin{center}
\begin{tikzpicture}
\begin{axis}[
  width=0.78\linewidth,
  height=7cm,
  axis lines=middle,
  xmin=0,
  xmax=4.2,
  ymin=-0.5,
  ymax=9.5,
  xlabel={$t,\ \text{с}$},
  ylabel={$v_x,\ \text{м}/\text{с}$},
  xtick={0,1,2,3,4},
  ytick={0,1,3,5,7,9},
  grid=both,
  minor tick num=0,
  unit vector ratio=1 1,
  clip=false
]
\addplot[
  domain=0:4,
  samples=200,
  very thick,
  color=accent
]
{1+2*x};

\addplot[
  only marks,
  mark=*,
  mark size=2.2pt,
  color=darkaccent
]
coordinates {(0,1) (1,3)};
\end{axis}
\end{tikzpicture}
\end{center}

\subsection*{Пример 2. Торможение, остановка и смена направления}

Пусть
\[
  v_x=5-t.
\]

Тогда
\[
  v_{0x}=5\,\text{м}/\text{с},
  \qquad
  a_x=-1\,\text{м}/\text{с}^2.
\]

До момента
\[
  t=5\,\text{с}
\]
скорость положительна. Тело движется вдоль \(+Ox\), а модуль скорости уменьшается.

При
\[
  t=5\,\text{с}
\]
получаем
\[
  v_x=0.
\]
Тело остановилось.

После этого
\[
  v_x<0.
\]
Тело движется вдоль \(-Ox\). При этом \(|v_x|\) растёт, поэтому тело разгоняется в отрицательном направлении.

\begin{center}
\begin{tikzpicture}
\begin{axis}[
  width=0.82\linewidth,
  height=7cm,
  axis lines=middle,
  xmin=0,
  xmax=9,
  ymin=-4,
  ymax=6,
  xlabel={$t,\ \text{с}$},
  ylabel={$v_x,\ \text{м}/\text{с}$},
  xtick={0,1,2,3,4,5,6,7,8},
  ytick={-3,-2,-1,0,1,2,3,4,5},
  grid=both,
  minor tick num=0,
  clip=false
]
\addplot[
  domain=0:8,
  samples=200,
  very thick,
  color=accent
]
{5-x};

\addplot[
  only marks,
  mark=*,
  mark size=2.2pt,
  color=darkaccent
]
coordinates {(0,5) (5,0) (8,-3)};
\end{axis}
\end{tikzpicture}
\end{center}

% ============================================================

\section*{7. Как находить ускорение по графику}
\sectionline

Выберите на прямой две удобные точки:
\[
  (t_1;v_1),
  \qquad
  (t_2;v_2).
\]

Затем найдите изменения:
\[
  \Delta t=t_2-t_1,
\]
\[
  \Delta v_x=v_2-v_1.
\]

Ускорение:
\[
  \boxed{
  a_x=
  \frac{\Delta v_x}{\Delta t}
  }.
\]

\key{Геометрический смысл:}
ускорение является угловым коэффициентом графика \(v_x(t)\).

Если прямая возрастает,
\[
  a_x>0.
\]

Если прямая убывает,
\[
  a_x<0.
\]

Если график горизонтален,
\[
  a_x=0,
\]
то есть скорость постоянна.

% ============================================================

\section*{8. Как читать движение по графику \(v_x(t)\)}
\sectionline

Удобно каждый участок анализировать по трём вопросам.

\begin{enumerate}
  \item Каков знак \(v_x\)?
  \item Увеличивается или уменьшается \(|v_x|\)?
  \item Есть ли точка, где \(v_x=0\)?
\end{enumerate}

Например, если график идёт от
\[
  v_x=5\,\text{м}/\text{с}
\]
к
\[
  v_x=0
\]
за \(5\) секунд, а затем уходит в область отрицательных скоростей, движение описывается так:

\begin{itemize}
  \item от \(0\) до \(5\,\text{с}\) тело движется вдоль \(+Ox\) и замедляется;

  \item при \(t=5\,\text{с}\) тело останавливается;

  \item после \(5\,\text{с}\) тело движется против \(Ox\) и разгоняется.
\end{itemize}

\key{Важно:}
отрицательная скорость означает направление движения против положительного направления оси \(Ox\). Она сама по себе ничего не говорит о разгоне или торможении.

% ============================================================

\section*{9. Типичные ошибки}
\sectionline

\begin{enumerate}
  \item Найти время встречи и записать его как ответ, хотя требовалась координата встречи.

  \item Потерять знак скорости или ускорения. Знак имеет физический смысл и показывает направление относительно оси \(Ox\).

  \item Делать вывод о разгоне только по знаку \(a_x\). Нужно учитывать также знак \(v_x\).

  \item Считать, что в выражении
  \[
    v_x=-3+t
  \]
  коэффициента перед \(t\) нет. На самом деле
  \[
    v_x=-3+1\cdot t,
  \]
  поэтому
  \[
    a_x=1\,\text{м}/\text{с}^2.
  \]

  \item Забывать единицы измерения:
  \[
    v_x:\ \text{м}/\text{с},
  \]
  \[
    a_x:\ \text{м}/\text{с}^2,
  \]
  \[
    t:\ \text{с},
  \]
  \[
    x:\ \text{м}.
  \]

  \item При построении графика пытаться заранее «вычислить \(t\)». При составлении таблицы значения независимой переменной \(t\) выбираются самостоятельно.

  \item Ошибаться при переносе слагаемых и работе с дробями.

  \item Смешивать скорость и модуль скорости. Например,
  \[
    -4<-2,
  \]
  но
  \[
    |-4|>|-2|.
  \]
  Поэтому тело со скоростью \(-4\,\text{м}/\text{с}\) движется быстрее по модулю, чем тело со скоростью \(-2\,\text{м}/\text{с}\).
\end{enumerate}

% ============================================================

\section*{10. Памятка перед решением}
\sectionline

\begin{itemize}
  \item Выпишите данные задачи.

  \item Переведите величины в СИ.

  \item Запишите общую физическую формулу.

  \item Подставьте известные числа.

  \item Неизвестную величину оставьте буквой.

  \item Решите полученное алгебраическое уравнение.

  \item Проверьте знак результата.

  \item Проверьте единицы измерения.

  \item Для графика \(v_x(t)\) помните:
  \[
    v_x=v_{0x}+a_xt.
  \]

  \item Свободный член --- это \(v_{0x}\).

  \item Коэффициент перед \(t\) --- это \(a_x\).

  \item Если график пересекает ось времени, в точке пересечения
  \[
    v_x=0.
  \]
  Это момент остановки и возможной смены направления.
\end{itemize}

% ============================================================
% ТРЕНИРОВОЧНЫЙ БЛОК
% ============================================================

\clearpage

\section*{Тренировочный блок. Путь Б}
\sectionline

На каждый день предусмотрено по 5 заданий.
Решения оформляйте по схеме:

\[
\text{формула}
\rightarrow
\text{подстановка}
\rightarrow
\text{вычисление}
\rightarrow
\text{ответ с единицами}.
\]

% ------------------------------------------------------------

\subsection*{13.09.2026}

\begin{enumerate}[label=\textbf{\arabic*.}]

  \item Скорость тела увеличилась с
  \(6\,\text{м}/\text{с}\)
  до
  \(18\,\text{м}/\text{с}\)
  за
  \(4\,\text{с}\).
  Найдите ускорение.

  \item Тело движется с начальной скоростью
  \(4\,\text{м}/\text{с}\)
  и постоянным ускорением
  \(3\,\text{м}/\text{с}^2\).
  Найдите скорость через
  \(5\,\text{с}\).

  \item Конечная скорость равна
  \(30\,\text{м}/\text{с}\),
  начальная скорость ---
  \(10\,\text{м}/\text{с}\),
  ускорение ---
  \(5\,\text{м}/\text{с}^2\).
  Найдите время разгона.

  \item Скорость задана формулой
  \[
    v_x=-2+4t.
  \]
  Найдите \(v_{0x}\), \(a_x\) и момент остановки.
  Укажите направление движения до остановки и после неё.

  \item Два тела движутся по законам
  \[
    x_1=18-2t,
  \]
  \[
    x_2=-6+4t.
  \]
  Найдите время и координату встречи.

\end{enumerate}

% ------------------------------------------------------------

\subsection*{14.09.2026}

\begin{enumerate}[label=\textbf{\arabic*.}]

  \item Скорость тела уменьшилась с
  \(25\,\text{м}/\text{с}\)
  до
  \(10\,\text{м}/\text{с}\)
  за
  \(5\,\text{с}\).
  Найдите ускорение.

  \item Тело имело начальную скорость
  \(-6\,\text{м}/\text{с}\)
  и ускорение
  \(2\,\text{м}/\text{с}^2\).
  Найдите скорость через
  \(4\,\text{с}\)
  и укажите направление движения в этот момент.

  \item Конечная скорость тела равна
  \(14\,\text{м}/\text{с}\),
  ускорение ---
  \(2\,\text{м}/\text{с}^2\),
  время движения ---
  \(6\,\text{с}\).
  Найдите начальную скорость.

  \item Скорость задана формулой
  \[
    v_x=7-t.
  \]
  Найдите начальную скорость, ускорение, момент остановки и опишите движение до остановки и после неё.

  \item На графике \(v_x(t)\) известны две точки:
  \[
    (1;6),
    \qquad
    (4;0).
  \]
  Найдите ускорение. Затем определите начальную скорость и запишите закон \(v_x(t)\).

\end{enumerate}

% ------------------------------------------------------------

\subsection*{15.09.2026}

\begin{enumerate}[label=\textbf{\arabic*.}]

  \item Скорость тела изменилась с
  \(-4\,\text{м}/\text{с}\)
  до
  \(8\,\text{м}/\text{с}\)
  за
  \(6\,\text{с}\).
  Найдите ускорение.

  \item Тело движется по закону
  \[
    v_x=3-2t.
  \]
  Найдите момент остановки и скорость через
  \(4\,\text{с}\).
  Опишите направление движения до остановки и после неё.

  \item За какое время тело с начальной скоростью
  \(5\,\text{м}/\text{с}\)
  при ускорении
  \(-1{,}5\,\text{м}/\text{с}^2\)
  достигнет скорости
  \(-4\,\text{м}/\text{с}\)?

  \item Через
  \(8\,\text{с}\)
  скорость тела стала равной
  \(-12\,\text{м}/\text{с}\).
  Ускорение равно
  \(-2\,\text{м}/\text{с}^2\).
  Найдите начальную скорость.

  \item Два тела движутся по законам
  \[
    x_1=-10+3t,
  \]
  \[
    x_2=14-t.
  \]
  Найдите время и координату встречи.
  Проверьте координату подстановкой в оба закона движения.

\end{enumerate}

% ============================================================
% ОТВЕТЫ
% ============================================================

\clearpage

\section*{Ответы к тренировочному блоку}
\sectionline

\begin{center}
\begin{tabularx}{\linewidth}{@{}p{25mm}p{11mm}X@{}}
\toprule
Дата & № & Ответ \\
\midrule

13.09.2026
& 1
& \(3\,\text{м}/\text{с}^2\) \\

& 2
& \(19\,\text{м}/\text{с}\) \\

& 3
& \(4\,\text{с}\) \\

& 4
& \(v_{0x}=-2\,\text{м}/\text{с}\);
\(a_x=4\,\text{м}/\text{с}^2\);
\(t_{\text{ост}}=0{,}5\,\text{с}\).
До остановки тело движется против \(Ox\) и замедляется; после остановки движется вдоль \(+Ox\) и разгоняется. \\

& 5
& \(t=4\,\text{с}\);
\(x=10\,\text{м}\) \\

\midrule

14.09.2026
& 1
& \(-3\,\text{м}/\text{с}^2\) \\

& 2
& \(2\,\text{м}/\text{с}\);
движение вдоль \(+Ox\) \\

& 3
& \(2\,\text{м}/\text{с}\) \\

& 4
& \(v_{0x}=7\,\text{м}/\text{с}\);
\(a_x=-1\,\text{м}/\text{с}^2\);
\(t_{\text{ост}}=7\,\text{с}\).
До остановки тело движется вдоль \(+Ox\) и замедляется; после остановки движется против \(Ox\) и разгоняется. \\

& 5
& \(a_x=-2\,\text{м}/\text{с}^2\);
\(v_{0x}=8\,\text{м}/\text{с}\);
\[
v_x=8-2t
\]
\\

\midrule

15.09.2026
& 1
& \(2\,\text{м}/\text{с}^2\) \\

& 2
& \(t_{\text{ост}}=1{,}5\,\text{с}\);
\(v_x(4)=-5\,\text{м}/\text{с}\).
До остановки тело движется вдоль \(+Ox\) и замедляется; после остановки движется против \(Ox\) и разгоняется. \\

& 3
& \(6\,\text{с}\) \\

& 4
& \(4\,\text{м}/\text{с}\) \\

& 5
& \(t=6\,\text{с}\);
\(x=8\,\text{м}\) \\

\bottomrule
\end{tabularx}
\end{center}

\end{document}